\section{Introduction}

Light is an electromagnetic wave whose electric field can exhibit a specific orientation, a property known as polarization. The polarization of light is a
fundamental aspect of optics and plays an important role in a wide range of scientific and technological applications. Examples include polarized sunglasses
that reduce glare, optical microscopy techniques that enhance image contrast, modern telecommunications systems that use polarization for signal transmission,
and emerging quantum technologies where polarization states can serve as carriers of information.
\\
The aim of this experiment is to investigate the fundamental properties of polarized light. Different methods for generating polarized light will be explored,
and the effects of various optical elements on its polarization state will be examined. Furthermore, techniques for measuring and analyzing polarization will 
be introduced. Through these investigations, a deeper understanding of the behavior of polarized light and its practical applications in optical systems will 
be developed.